\section{Discussion}
\label{sec:discussion}

\paragraph{What does healing restore?}
In these observed runs, the combined evidence argues against treating recovery
as a single scalar. PPL
measures how well the healed model fits held-out next-token events; it improves
substantially in every inherited-prefix arm. That improvement is useful---the
models recover fluent continuation and several commonsense tasks---but it does
not summarize answer selection or closed-book recall by itself. The endpoint
scatter makes the distinction concrete: models with similar PPL can occupy
meaningfully different MMLU regimes, and random initialization can outrank a
frozen inherited front in PPL while underperforming it on answer letters. We
therefore interpret healing as partial recovery of distributional modeling plus
a task-dependent remainder in this case study, not as a binary return of the
original model or a matched-PPL law.

\paragraph{Interface sensitivity is substantial but insufficient.}
Complete-option MMLU recovers more smoothly than answer-letter MMLU, showing
substantial protocol sensitivity that is consistent with an answer-symbol
component. The protocols also change the prompt, candidate continuation,
tokenization, and normalization, so they do not isolate mapping loss. The
random-init fluency floor and the independent PopQA, TriviaQA,
and NQ-open gaps argue against a formatting-only account. This distinction matters
for evaluation: a single multiple-choice protocol can overstate damage when its
interface is brittle, whereas a content-likelihood protocol can overstate
recovery when generic fluency makes plausible options easy to score. Agreement
across interfaces is more informative than designating either as ground truth.

\paragraph{Different constructions yield different endpoints.}
At the same nominal 200k step budget, the keep14--ShortGPT comparison shows that
these two half-depth constructions are not one operating point. ShortGPT yields a much stronger
likelihood--capability endpoint than the tested prefix-plus-fresh-tail recipe.
The result does not identify whether the advantage comes from the
original final block, two additional inherited blocks, non-contiguous selection,
or avoiding fresh layers. It shows that conclusions drawn from one prefix
construction should not be generalized to every 16-layer construction;
isolating the responsible factor requires additional matched structures.

\paragraph{Which behaviors remain difficult?}
The residual deficit spans factual recall (PopQA, TriviaQA, and NQ-open), broad
academic knowledge (all four MMLU groups), and some answer-interface behavior.
In contrast, several continuation-style commonsense tasks recover more readily.
This taxonomy is behavioral rather than anatomical. Knowledge-neuron and causal
tracing studies motivate asking whether particular feed-forward computations
support factual recall~\citep{dai2022knowledge,meng2022locating}, but neither our
subject groups nor ShortGPT's retention of block 31 identifies the missing
substrate. Structure-isolation and causal interventions are required before
making that claim.

\paragraph{A reporting proposal from this case study.}
The measurements motivate reporting likelihood, target capability, interface,
recovery compute, and exact structural construction together. An intact model
continued on the same corpus can bound drift only over its observed horizon;
same-shape operating points can expose alternative explanations only to the
extent that learning rate and trainable modules are matched.

\begin{center}
\fbox{\begin{minipage}{0.92\columnwidth}
\small\textbf{For a more interpretable recovery claim:}
(1) evaluate an independent target capability;
(2) test a second scoring or generation interface;
(3) state the horizon and matching of any continued-full-model control; and
(4) report the exact structure and healing budget.
\end{minipage}}
\end{center}

This proposal does not replace efficiency or endpoint-quality comparisons and
is not validated across methods or model families. Its purpose is to prevent a
low PPL alone from being interpreted as evidence that a target behavior has
returned.
